\documentclass[11pt]{article}
\usepackage[margin=1in]{geometry}
\usepackage{amsmath,amssymb,amsthm}
\usepackage[colorlinks=true,linkcolor=blue,citecolor=blue,urlcolor=blue]{hyperref}
\newtheorem{theorem}{Theorem}
\newtheorem{lemma}{Lemma}
\newtheorem{corollary}{Corollary}
\newcommand{\one}{\mathbf1}
\newcommand{\Prob}{\mathbb P}
\newcommand{\E}{\mathbb E}
\newcommand{\Fp}{\mathbb F_p}
\newcommand{\supp}{\operatorname{supp}}
\title{Boolean rigidity and polynomial kernel enumeration in the regular SBM}
\author{Research session note: self-contained type-counting argument}
\date{November 15, 2026}
\begin{document}
\maketitle
\begin{abstract}
For fixed internal degree $a>b>0$ with $a\ge3$, we prove that the planted signs are the only sign vectors at adjacency eigenvalue $a-b$, with polynomially small failure probability, at every admissible block size. Exact configuration-model type counting gives a strictly negative macroscopic entropy exponent; a small-support estimate handles near-planted alternatives. The same count over a fixed finite field bounds the expected kernel size polynomially, allowing polynomial-time enumeration of its sign vectors. This proves exact recovery without requiring one-dimensionality of the real eigenspace. All statements here use the original independent random cross graph, not an arbitrary prescribed cross graph. No novelty claim is made without further literature review.
\end{abstract}

Fix $a>b>0$, put $D=a+b$, $\delta=a-b$, and let both blocks have size $N$. All limits follow admissible sizes $aN$ even. The two internal graphs are independent uniform simple $a$-regular graphs; the cross graph is independent uniform simple $b$-regular bipartite. Denote the adjacency matrix by $A$ and planted signs by $\sigma=(\one,-\one)$.

\begin{theorem}[Boolean rigidity]\label{thm:rigidity}
If $a\ge3$, then
\[
 \Prob\{\exists\tau\in\{-1,1\}^{2N}\setminus\{\sigma,-\sigma\}:
 A\tau=\delta\tau\}
 =O_{a,b}\bigl(N^{-(D-2)(\delta+1)}\bigr).
\]
In particular this includes $a=4$ and odd $N$.
\end{theorem}

Every sign eigenvector at $\delta$ is balanced, because
$(D-\delta)\one^T\tau=0$ over $\mathbb R$ and $D-\delta=2b>0$.
Thus an alternative flips $k$ vertices in each planted block; replacing it by its negative allows $1\le k\le N/2$.

\section{Exact configuration-model type count}
Start with independent internal configuration pairings and a uniform cross-stub bijection. Give each vertex $a$ ordered internal stubs and $b$ ordered cross stubs. For an alphabet $\mathcal S$, a star profile records its root value, its $a$ internal neighbor values, and its $b$ cross neighbor values. Retain only profiles obeying
\begin{equation}\label{eq:local}
 \sum_{i=1}^a x_i+\sum_{j=1}^b y_j=\delta x_0.
\end{equation}
Use real equality for $\mathcal S=\{-1,1\}$ and finite-field equality for $\mathcal S=\Fp$.

Let $\pi_+,\pi_-$ be empirical star-profile laws in the two blocks. Their root marginals are $\mu_+,\mu_-$. Let $m^+_{st}$ count plus-block internal stubs with root value $s$ and desired neighbor value $t$, and similarly $m^-$. Let $e_{st}$ count plus-block cross stubs with root value $s$ and desired minus-neighbor value $t$. Feasible pairings require $m^\pm$ symmetric with even diagonal, and the opposite cross table equal to $e^T$. Put
\[
 Q_\pm=\frac{m^\pm}{aN},\qquad R=\frac e{bN}.
\]
The marginals of $Q_\pm$ are $\mu_\pm$ in both coordinates; those of $R$ are $\mu_+,\mu_-$.

For assigned labeled star profiles, the number of compatible internal pairings for a table $m$ is
\[
 W(m)=\prod_{s}(m_{ss}-1)!!\prod_{s<t}m_{st}!,
\]
and the number of cross bijections is $\prod_{s,t}e_{st}!$. Write $n^\pm_\omega=N\pi_\pm(\omega)$ for profile counts. The exact expected number of alphabet-valued eigenvectors is the sum over all feasible integer types of
\begin{equation}\label{eq:exact}
 \frac{N!}{\prod_\omega n^+_\omega!}
 \frac{N!}{\prod_\omega n^-_\omega!}
 \frac{W(m^+)W(m^-)\prod_{s,t}e_{st}!}
 {((aN-1)!!)^2(bN)!}.
\end{equation}
Indeed, choose the profiles at labeled vertices and then pair the stubs by their root and desired-neighbor classes. Each actual eigenvector and configuration determines those labels uniquely. Root-label assignments are already included in the two multinomial factors; there is no omitted root-entropy term.

For fixed alphabet and degrees, there are polynomially many types. Uniform Stirling estimates, including zero counts, give the logarithm of each term in \eqref{eq:exact} as
\begin{equation}\label{eq:Phi}
 N\Phi+O(\log(N+1)),\qquad
 \Phi=H(\pi_+)+H(\pi_-)-\frac a2\{H(Q_+)+H(Q_-)\}-bH(R),
\end{equation}
where entropies use natural logarithms. The error constant depends only on the fixed alphabet and degrees.

Independently symmetrize the internal stub positions and cross stub positions in each profile law. This preserves \eqref{eq:local}, all aggregated edge tables, and root marginals, while increasing profile entropy. It therefore suffices to upper-bound $\Phi$ for such symmetric laws.

\section{The entropy inequality}
For a symmetric plus-star law write the variables as
$(X;Y_1,\ldots,Y_a;Z_1,\ldots,Z_b)$. Because \eqref{eq:local} determines $Y_a$ from all the other entries,
\begin{align*}
 H(\pi_+)
 &=H(X)+H(Y_1,\ldots,Y_{a-1},Z_1,\ldots,Z_b\mid X)\\
 &\le H(\mu_+)+(a-1)H(Y_1\mid X)+bH(Z_1\mid X).
\end{align*}
Let $J_+\ge0$ be the conditional subadditivity defect in this inequality; define $J_-$ analogously. Put
\[
 I_+=2H(\mu_+)-H(Q_+),\quad
 I_-=2H(\mu_-)-H(Q_-),\quad
 I_\times=H(\mu_+)+H(\mu_-)-H(R).
\]
These are mutual informations, hence nonnegative. Substitution into \eqref{eq:Phi} gives the exact identity for symmetrized laws
\begin{equation}\label{eq:MI}
 -\Phi=\left(\frac a2-1\right)(I_++I_-)+bI_\times+J_++J_-.
\end{equation}
Thus $\Phi\le0$ for every alphabet whenever $a\ge2$.

For sign profiles with flip fraction $\alpha=k/N$,
\[
 \mu_+(1)=1-\alpha,\qquad \mu_-(1)=\alpha.
\]
When $a\ge3$, equality $\Phi=0$ would force all three mutual informations and both defects to vanish. Conditional on $X$, the retained $a-1$ internal values and $b$ cross values are then independent, with laws $\mu_+$ and $\mu_-$. The eliminated $Y_a$ also has law $\mu_+$ conditional on $X$, by $I_+=0$. Conditional variances in
$Y_a=\delta X-\sum_{i<a}Y_i-\sum_jZ_j$ would give
\[
 \operatorname{Var}(\mu_+)=(a-1)\operatorname{Var}(\mu_+)+b\operatorname{Var}(\mu_-).
\]
Hence both variances vanish. The constant laws must be opposite by \eqref{eq:local}, so they are the planted assignment or its global negative. This already implies a strictly negative entropy maximum on every compact macroscopic overlap range.

An explicit bound is also available. Set $L=-\Phi$, $v=4\alpha(1-\alpha)$, and $K=\delta^2+D-2$. Let $P$ be the joint law of the retained plus-star entries $(X,Y_1,\ldots,Y_{a-1},Z_1,\ldots,Z_b)$, and let $T$ be the product law with the same one-coordinate marginals. Then
\[
 D_{\rm KL}(P\Vert T)=J_++(a-1)I_++bI_\times\le4L,
\]
using $(a-1)/(a/2-1)\le4$ for $a\ge3$. Thus Pinsker's inequality gives $\|P-T\|_{\rm TV}\le\sqrt{2L}$. For
$U=\delta X-\sum_{i<a}Y_i-\sum_jZ_j$, we have $U^2=1$ under $P$, whereas independence and the opposite root means give
\[
 \E_T U^2-1=Kv.
\]
Since $|U|\le2a-1$, comparison of the bounded function $U^2$ under the two laws yields
\begin{equation}\label{eq:gap}
 \Phi\le-\frac{8(\delta^2+D-2)^2}{(2a-1)^4}
 \alpha^2(1-\alpha)^2.
\end{equation}
For nonsymmetric laws the same upper bound follows by symmetrizing first.

Consequently, for every fixed $\eta>0$, the expected number of sign eigenvectors with $\eta N\le k\le N/2$ is at most $\exp(-c_{a,b,\eta}N)$ for all sufficiently large $N$, by \eqref{eq:exact}--\eqref{eq:gap} and the polynomial number of types.

\section{Near-planted alternatives}
\begin{lemma}[Small-support bound]\label{lem:support}
For some fixed $\varepsilon=\varepsilon(a,b)>0$,
\[
 \Prob\{\exists x\ne0:Ax=\delta x,\ |\supp x|\le\varepsilon N\}
 =O_{a,b}\bigl(N^{-(D-2)(\delta+1)}\bigr).
\]
\end{lemma}
\begin{proof}
In a simple $D$-regular graph, a $\delta$-eigenvector has zero coordinate sum and both signs. At its maximum positive coordinate there must be at least $\delta$ positive neighbors; apply the same observation to its negative. Thus its support size $q$ is at least $q_0=2(\delta+1)$.

For a fixed possible support $S$ of size $q$, every exterior vertex adjacent to $S$ must have at least two incident edges from $S$, since otherwise the eigenvector equation at that zero coordinate cannot vanish. Expose the configuration pairings incident to $S$ adaptively. A collision means pairing to $S$ or to an exterior vertex already seen. If $i$ exposed edges lie within $S$, $j$ go outside, and $t$ distinct exterior vertices are reached, the number of collisions is $i+j-t$. Necessarily $t\le j/2$ and $Dq=2i+j$, so at least $Dq/2$ collisions occur in at most $Dq$ exposures.

For $q\le\varepsilon N$ with $\varepsilon$ small enough, each conditional collision probability is at most $Cq/N$: at most $O_D(q)$ stubs lie at previously reached vertices, while every relevant internal or cross pool retains a positive fraction of its $\Theta(N)$ stubs. A union bound over collision positions and then over $\binom{2N}q$ supports gives
\[
 \Prob(\exists\text{ such support of size }q)
 \le\left[C'\left(\frac qN\right)^{(D-2)/2}\right]^q.
\]
Choose $\varepsilon$ small enough that successive terms decrease geometrically. Summing from $q_0$ to $\varepsilon N$ gives $O(N^{-(D-2)(\delta+1)})$.

The support lower bound was used only for the desired simple-graph event. The exploration bound was performed in the configuration model. Conditioning its two internal factors and bipartite factor on simplicity costs a bounded factor at fixed degrees and yields the requested independent uniform-simple model.
\end{proof}

For an alternative sign vector near $\sigma$, take $x=\sigma-\tau$. Its support has size $2k$, so Lemma~\ref{lem:support} handles $k\le\varepsilon N/2$. The preceding macroscopic argument handles the remaining $k$. Fixed-degree configuration simplicity has probability bounded below, so the macroscopic exponential bound also survives conditioning. This proves Theorem~\ref{thm:rigidity}.

\section{Polynomial finite-field moment and enumeration}
\begin{theorem}\label{thm:polymoment}
Fix $a\ge2$, $b>0$, and an odd prime $p$. For some finite constants $C,\kappa$ depending only on $a,b,p$,
\[
 \E\bigl|\ker_{\Fp}(A-\delta I)\bigr|\le C N^\kappa
\]
at all sufficiently large admissible sizes, in the original uniform-simple RSBM.
\end{theorem}
\begin{proof}
Use alphabet $\Fp$ in \eqref{eq:exact}. There are at most $p^{D+1}$ star profiles in each block and polynomially many integer types. The uniform Stirling prefactor is polynomial in $N$, with a fixed exponent. Equation~\eqref{eq:MI} gives $\Phi\le0$ for every feasible type when $a\ge2$. Summing \eqref{eq:exact} therefore gives the polynomial bound in the configuration model. Bounded-cost conditioning on simplicity proves the result as stated. No distributional limit or uniform growing-prime estimate is needed.
\end{proof}

Let $r=\dim_{\Fp}\ker(A-\delta I)$. Markov's inequality gives, for any desired fixed $t>0$, a fixed computable constant $c=c(a,b,p,t)$ such that
\[
 \Prob(r>c\log N)=O(N^{-t}).
\]
The conservative constants can be read from the finite type count and Stirling bounds; no optimization is needed for polynomial time.

\begin{corollary}[Recovery for internal degree at least three]
For every fixed $a>b>0$ with $a\ge3$, there is a polynomial-time algorithm recovering the planted partition with high probability, at every admissible size, without an outlier-gap assumption.
\end{corollary}
\begin{proof}
Fix an odd prime, for example $p=3$. Compute a basis matrix for the finite-field kernel. If $r>c\log N$, report failure. Otherwise choose $r$ coordinate rows giving an invertible $r\times r$ submatrix. Enumerate all $2^r$ assignments of signs to those coordinates, extend each uniquely within the kernel, retain vectors whose every coordinate is $\pm1$, and lift those signs to integers. Verify the exact integer equation $A\tau=\delta\tau$ for each candidate. This verification removes any extra modular solutions, even when $p$ is smaller than the degree.

The planted signs are always enumerated. On the event of Theorem~\ref{thm:rigidity}, the only surviving vectors are $\pm\sigma$. The cutoff event has polynomially small probability by Theorem~\ref{thm:polymoment}. With the cutoff the worst-case running time is polynomial, since $2^r\le N^{c\log2}$.
\end{proof}

For $(a,b)=(2,1)$, the separate exact counting proof in the companion low-degree note establishes Boolean rigidity with ambiguity probability $\Theta(N^{-2})$. Theorem~\ref{thm:polymoment} and the same enumeration then apply to that pair as well. This last statement uses the companion proof; the internal-degree-at-least-three rigidity theorem above does not silently cover $a=2$.

These recovery conclusions do not prove that the real planted eigenspace is one-dimensional in the previously uncovered cases. They also do not extend the Boolean-rigidity or polynomial-moment claims to an adversarial prescribed cross graph: the exact count uses the independent random bipartite factor.
\end{document}
